\documentclass{article}
\usepackage{mathnotes}

\title{Markov Chains}
\description{Stochastic processes with transition matrices, ergodic properties, and steady-state probability calculations using linear algebra methods.}

\begin{document}

\section{Markov Chains}

\begin{definition}[Markov Chain]
If we have a system with N states and for each state, a set of N \@{probabilities} from transitioning from that state to another state (including itself) we end up with an $N \times N$ \@{matrix} $P$ where the \@{entry} $P_{jk}, j,k \in \{1, \dots, N\}$ represents the probability of transitioning from state $j$ to state $k$.

Starting with some state, we can transition to another state, then another, with the probability of picking the next state dependent only on the existing state and the matrix $P.$ We call this sequence of random events a \textbf{Markov Chain}.

	More formally, the \@{sequence} or \@{random-variables}  $X_1, X_2, \dots, X_n$ is called a Markov Chain if

\[ P(X_{n+1} = k|X_n = j, X_{n-1} = j_{n-1}, \dots, X_0 = j_0) = P(X_{n+1} = k | X_n = j). \]

We write

\[ p_{jk} = P(X_{n+1} = k | X_n = j), \]

and note that

\[ p_{jk} \geq 0 \]

and

\[ \sum_{k=1}^{N} p_{jk} =  1 \]

for $j = 1,2, \dots, N.$ We call the $p_{jk}$ values the transition probabilities of the Markov chain.

	Alternatively, we can say the \@{sequence} or \@{random-variables}  $X, Y, Z$ is said to form a Markov Chain (denoted $X \to Y \to Z$) if and only if $X$ and $Z$ are \@{conditionally independent} given $Y.$ Specifically, $X, Y$ and $Z$ form a Markov Chain $X \to Y \to Z$ if the \@{joint probability mass function} can be written as

	\[ p(x, y, z) = p(x) p(y|x) p(z|y). \]

\end{definition}

\begin{theorem}[Conditional Independence Given Intermediate Step]
	\@{Random Variables} $X, Y, Z$ form a \@{markov chain} $X \to Y \to Z$ if and only if $X$ and $Z$ are \@{conditionally independent} given $Y.$
\end{theorem}
\begin{proof}
	Suppose $X$ and $Z$ are \@{conditionally independent} given $Y.$ Then, using \@{bayes theorem} we have

	\bal
	p(x, z | y) & = p(x|y)p(z|y) \\
	            & = \frac{p(x)p(y|x) p(z|y)}{p(y)}
	\eal

	Then,

	\bal
	p(x,y,z) & = p(y)p(x|y)p(z|y) \\
		 & =  p(y) \frac{p(x)p(y|x) p(z|y)}{p(y)} \\
		 & = p(x)p(y|x)p(z|y).
	\eal

	Now assume we have Markovicity, i.e $p(x,y,z) = p(x)p(y|x)p(z|y).$ Again, using Bayes:

	\bal
	p(x,y,z) & = p(x)p(y|x)p(z|y) \\
	         & = \frac{p(x) p(x|y) p(y) p(z|y)}{p(x)} \\
	         & = p(x|y) p(y) p(z|y).
	\eal

	So, for a fixed $y$ we have
	\[ \frac{p(x,y,z)}{p(y)} = p(x|y)p(z|y). \]
\end{proof}

\begin{theorem}[Reversibility of Markov Chains]
	If \@{Random Variables} $X, Y, Z$ form a \@{markov chain} $X \to Y \to Z,$ then $Z \to Y \to X$ is also a Markov Chain (hence it is sometimes written as $X \leftrightarrow Y \leftrightarrow Z.)$
\end{theorem}
\begin{proof}
Using \@{bayes theorem},

\bal
	p(x, z | y) & = p(x|y)p(z|y) \\
		    & = \frac{p(x|y)p(y|z)p(z)}{p(y)}.
\eal
	Then,

\bal
	p(x,y,z) & = p(y)p(x|y)p(z|y) \\
		 & = p(y) \frac{p(x|y)p(y|z)p(z)}{p(y)}. \\
		 & = p(z)p(y|z)p(x|y).
\eal
\end{proof}

\begin{theorem}
Given \@{Random Variables} $X, Y, Z$ if $Z = f(Y),$ then $X \to Y\ \to Z.$
\end{theorem}
\begin{proof}
If $Z = f(Y),$ then $Z$ is independent of $X,$ so

	\[ p(y,z | x) = p(y|x,z) = p(y|x)p(z|y), \]

	and then

	\[ p(x,y,z) = p(x)p(y|x)p(z|y). \]
\end{proof}


\detach

\begin{note}
Note that $p_{jk}$ is the \@{probability} of going from state $j$ to state $k$ in one step. We use the notation

\[ p_{jk}^{(m)} \]

to represent the probability of going from state $j$ to state $k$ in $m$ steps.
\end{note}

\begin{theorem}
$p_{jk}^{(m)} = P_{jk}^m.$
\end{theorem}

\begin{proof}
Due to the magic of linear algebra, if we have a row \@{vector} $\vec{x}$ representing the \@{probability} of starting in each state, we can compute the probability of ending up in any given state after $m$ steps as

\[ \vec{x} \cdot P^m. \]
\end{proof}

\subsection{Ergodic Markov Chains}

\begin{definition}[Ergodic]
If for some positive integer $m$ we have that $p_{jk}^{(m)} > 0$ for all $j,k = 1,2,\dots,N,$ then the \@{Markov-chain} is said to be \textbf{ergodic.}
\end{definition}

\begin{theorem}[Stationary Distribution of an Ergodic Chain]\label{stationary-distribution}
When the \@{markov-chain} is \@{ergodic}, the \@[limit]{limit-sequence}

\[ \pi_k = \lim_{n \to \infty} p_{jk}^{(n)} \]

exists and

\[ \sum_{k=1}^N \pi_k = 1. \]
\end{theorem}

\begin{note}
Here, $\vec{\pi} = \langle \pi_1, \pi_2, \dots, \pi_N \rangle$ and the \@{entries} in it represent the long term \@{probabilities} of being in any given state, or equivalently, the portion of time spent in any given state.

Note that $\vec{\pi} \cdot P = \vec{\pi},$ which when combined with the fact that the entries of $\vec{\pi}$ sum to 1 lets us solve a linear system of equations to find the entries of $\vec{\pi}.$
\end{note}

\begin{definition}[Probability Vector]
	A \@{vector} whose \@{entries} are \@{non-negative} and sum to 1 is called a \textbf{probability vector}.
\end{definition}

\begin{theorem}\label{probability-vectors-form-a-convex-compact-set}
The set of \@{probability vectors} in $\reals^n$ is a \@{convex-set}, and it is \@{compact}.


	\begin{proof}
		Let $E$ be the set of \@{probability vectors} in $\reals^n.$ That is,

		\[ E = \left \{ \vec{p} \in \reals^n : \vec{p}_i \geq 0, \sum_{i=1}^n \vec{p}_i = 1. \right \} \]

		(Note: $E$ is a $n-1$ dimensional \@{simplex}).

		Suppose $\vec{X}, \vec{Y} \in E.$ Then

		\[ \sum_{i=1}^n \vec{X}_i = 1, \quad \sum_{i=1}^n \vec{Y}_i = 1 \].

		Then, let $\lambda \in [0,1].$ Now,

		\[ \sum_{i=1}^n \lambda  \vec{X}_i = \lambda,  \quad \sum_{i=1}^n (1 - \lambda) \vec{Y}_i = 1 - \lambda, \]

		\bal
		\lambda \vec{X} + (1 - \lambda) \vec{Y} & = \sum_{i=1}^n \left ( \lambda  \vec{X}_i + (1 - \lambda) \vec{Y}_i \right ) \\
		& = \lambda + 1 - \lambda \\
		& = 1.
		\eal

		Since either $\lambda$ or $1 - \lambda$ are non-negative, all entries of this \@{convex combination} are non-negative, so the
		combination is again in $E$ and is therefore a \@{probability-vector}.

		We can use \@{heine-borel} to show that $E$ is \@{compact} by showing it is \@{closed} and \@{bounded}.

		Noe that for each $\vec{p} \in E,$ and each \@{coordinate} $i,$

		\[ \{ \vec{p} \in \reals^n : \vec{p}_i \geq 0 \} \]

		is a \@{closed} half-space. Now, note that the map

		\[ f : \reals^n \to \reals, \quad f(\vec{p}) = \sum_{i=1}^n \vec{p}_i \]

		is \@{continuous}, the \@{pre-image} of $1$ under $f$ is

		\[ \left \{ \vec{p} \in \reals^n : \sum_{i=1}^n \vec{p}_i = 1. \right \} = f^{-1}(\{1\}). \]

		Because ${1}$ is \@{closed} in $\reals,$ by \@{mapping-continuous-iff-inverse-images-of-closed-sets-are-closed} we have that

                \[ \left \{ \vec{p} \in \reals^n : \sum_{i=1}^n \vec{p}_i = 1. \right \} \]

		is \@{closed.}

		Now,

		\[ E = \left ( \bigcap_{i=1}^n \{ \vec{p} : \vec{p}_i \geq 0 \} \right ) \cap \left \{ \vec{p} : \sum_{i=1}^n \vec{p}_i = 1 \right \} \]

		is an \@{intersection} of \@{closed} \@{sets}, \@[which is also closed]{union-and-intersection-of-open-and-closed-sets}.

		$E$ is obviously bounded - the sum of all coordinates in $\vec{p} \in E$ is always $1,$ so $E \subset$ the \@{closed} \@{ball} of radius $1$ centered at the origin.

		Therefore, $E$ is closed and bounded, and is compact.


                TODO: we can generalize this to the fact that \@{simplexes} are \@{convex} and \@{compact}.
	\end{proof}

\end{theorem}

\begin{definition}[Stochastic Matrix]\synonyms{"probability matrix", "transition matrix", "Markov matrix", "right stochastic matrix", "row stochastic matrix"}
A \textbf{stochastic matrix} (actually, a right stochastic matrix) is a \@{square-matrix} $A = (a_{ij}) \in \reals^{n \times n}$ with \@{non-negative} \@{entries} whose rows each sum to one:

\[ a_{ij} \geq 0 \quad \text{for all } i,j, \qquad \sum_{j} a_{ij} = 1 \quad \text{for all } i. \]

Equivalently, $A\mathbf{1} = \mathbf{1}$ with $A \geq 0$ entrywise, where $\mathbf{1}$ is the all-ones \@{vector}.
\end{definition}

\begin{definition}[Left Stochastic Matrix]\synonyms{"column stochastic matrix"}
A \textbf{left stochastic matrix} is a \@{square-matrix} $A = (a_{ij}) \in \reals^{n \times n}$ with \@{non-negative} \@{entries} whose columns each sum to one:
\[ a_{ij} \geq 0 \quad \text{for all } i,j, \qquad \sum_{i} a_{ij} = 1 \quad \text{for all } j. \]
Equivalently, $\mathbf{1}^\top A = \mathbf{1}^\top$ with $A \geq 0$ entrywise, where $\mathbf{1}$ is the all-ones \@{vector}.
\end{definition}

\begin{note}
A \@{stochastic-matrix} can be used to describe the transitions of a \@{markov-chain}.
\end{note}

\begin{definition}[Doubly Stochastic Matrix]\synonyms{"bistochastic matrix"}
If a \@{matrix} is both a \@{row-stochastic-matrix} and a \@{column-stochastic-matrix}, it is said to be a \textbf{doubly stochastic matrix.} That is, both its rows and its columns sum to 1.
\end{definition}

\begin{note}
A \@{doubly-stochastic-matrix} is a \@{convex-combination} of \@{permutation-matrices}.
\end{note}

A \@{stochastic-matrix} maps a \@{probability vector} to another probability vector — but only for the matching multiplication convention. A \@{row-stochastic-matrix} preserves the sum acting on the right, $\vec{x}^\top A$; a \@{column-stochastic-matrix} preserves it acting on the left, $A\vec{x}$; and a \@{doubly-stochastic-matrix} preserves it either way. The demo below lets you apply each type to a probability vector, build doubly-stochastic matrices as \@{convex-combinations} of permutation matrices, and iterate the map toward its stationary distribution.

\includedemo{stochastic-matrix}

\end{document}
